\message{ !name(voorstel-inhoud.tex)}
\message{ !name(voorstel-inhoud.tex) !offset(249) }
De tweede fase richt zich op het identificeren van noden en vereisten van beide doelgroepen. Voor onderzoekers gaat het onder meer om transparantie van berekeningen, experimenteerruimte en begrijpelijke documentatie; voor ontwikkelaars om onderhoudbaarheid, stabiliteit en controle over data en interfaces. Deze noden worden gestructureerd geformuleerd als functionele en niet-functionele requirements, die richtinggevend zijn voor het verdere onderzoek. Hierbij wordt expliciet rekening gehouden met organisatorische aspecten, zoals taakverdeling en communicatiestromen, in lijn met inzichten uit onder meer \text
\message{ !name(voorstel-inhoud.tex) !offset(-4) }
